%============================================================
%  Chapter 6 – Backend Implementation
%============================================================
\chapter{Backend Implementation}

\section{Overview}

The \projectname{} backend is a Node.js/Express REST API written entirely in TypeScript. The entry point is \texttt{src/server.ts}, which bootstraps the HTTP server on a configurable port. The Express application is assembled in \texttt{src/app.ts}, where all global middleware and route groups are registered.

\section{API Route Groups}

The API is versioned under the \texttt{/api} prefix. Table~\ref{tab:routes} summarises all registered route groups.

\begin{table}[H]
\centering
\caption{API Route Groups}
\label{tab:routes}
\begin{tabularx}{\textwidth}{llX}
\toprule
\textbf{Prefix} & \textbf{File} & \textbf{Responsibility} \\
\midrule
\texttt{/api/auth}        & \texttt{auth.routes.ts}        & Registration, OTP, login, logout, OAuth, password reset. \\
\texttt{/api/projects}    & \texttt{project.routes.ts}     & CRUD for project listings; status transitions. \\
\texttt{/api/freelancers} & \texttt{freelancer.routes.ts}  & Freelancer onboarding, profile management, gigs. \\
\texttt{/api/proposals}   & \texttt{proposal.routes.ts}    & Submit, withdraw, shortlist, accept, reject proposals. \\
\texttt{/api/browse}      & \texttt{browse.routes.ts}      & Search/filter freelancers and projects with scoring. \\
\texttt{/api/track}       & \texttt{tracking.routes.ts}    & Record browse interactions, save/unsave projects. \\
\texttt{/api/tokens}      & \texttt{token.routes.ts}       & Token balance, history, admin grant/deduct. \\
\texttt{/api/categories}  & \texttt{category.routes.ts}    & Read-only category and subcategory lists. \\
\texttt{/api/metadata}    & \texttt{metadata.routes.ts}    & Languages and locations reference data. \\
\texttt{/api/ai}          & \texttt{ai.routes.ts}          & AI project scoping assistant endpoint. \\
\bottomrule
\end{tabularx}
\end{table}

\section{Authentication Controller}

The \texttt{auth.controller.ts} file implements the following endpoints:

\begin{description}
  \item[\texttt{POST /api/auth/signup}] Accepts \texttt{\{name, email, password\}}. Hashes the password with bcrypt, creates a \texttt{User} document with \texttt{isEmailVerified: false}, generates a 6-digit OTP, stores it in Redis with a 10-minute TTL, and sends it via Nodemailer.

  \item[\texttt{POST /api/auth/verify-otp}] Accepts \texttt{\{email, otp\}}. Looks up the OTP in Redis, compares it, sets \texttt{isEmailVerified: true}, deletes the OTP key, and returns JWT tokens.

  \item[\texttt{POST /api/auth/login}] Validates credentials, checks for ban status, generates an access token and refresh token, stores the refresh token as an HttpOnly cookie in a \texttt{Session} document.

  \item[\texttt{POST /api/auth/refresh}] Reads the refresh token from the HttpOnly cookie, validates it against the database, rotates the token (delete old, issue new), and returns a new access token.

  \item[\texttt{POST /api/auth/logout}] Deletes the session document and clears the cookie.

  \item[\texttt{POST /api/auth/forgot-password}] Generates a UUID-based reset token stored in \texttt{PasswordResetToken}, sends an email with a link containing the token. Token expires in 1 hour.

  \item[\texttt{POST /api/auth/reset-password}] Validates the reset token, hashes and updates the password, marks the token as used.

  \item[\texttt{GET /api/auth/google}] Redirects to Google's OAuth consent screen via Passport.js.

  \item[\texttt{GET /api/auth/google/callback}] Handles the OAuth callback. Uses Passport's Google strategy to find or create a \texttt{User}, then issues JWT tokens and redirects to the frontend.
\end{description}

\section{Freelancer Onboarding Controller}

\texttt{freelancer-onboarding.controller.ts} handles the multi-step onboarding flow. Key design points:

\begin{itemize}
  \item All onboarding steps are \texttt{PATCH} requests to step-specific endpoints (e.g., \texttt{/api/freelancers/onboarding/basic-info}, \texttt{/api/freelancers/onboarding/skills}).
  \item After each step, the \texttt{profileCompletion} score is recalculated using the \texttt{calculateProfileCompletion()} utility, which evaluates the number of filled mandatory fields as a fraction of total mandatory fields, multiplied by 100.
  \item A ``Save Draft'' mode accepts partial data without triggering mandatory field validation.
  \item File uploads (ID document, portfolio images) are handled by Multer in memory storage mode, then streamed to Cloudinary via \texttt{streamifier}.
\end{itemize}

\section{Project Controller}

\texttt{project.controller.ts} manages the full project lifecycle:

\begin{lstlisting}[language=TypeScript, caption={Project Routes Example}]
router.post('/', protect, requireRole('CLIENT'), createProject);
router.get('/', getAllProjects);
router.get('/:id', getProjectById);
router.patch('/:id', protect, requireRole('CLIENT'), updateProject);
router.delete('/:id', protect, requireRole('CLIENT'), deleteProject);
router.patch('/:id/publish', protect, requireRole('CLIENT'), publishProject);
\end{lstlisting}

The \texttt{createProject} handler accepts the full project specification, creates a \texttt{Project} document in \texttt{DRAFT} status, and optionally links \texttt{ProjectSkill} join documents for each required skill. The \texttt{publishProject} handler transitions status from \texttt{DRAFT} to \texttt{OPEN}.

\section{Proposal Controller}

The proposal workflow enforces the SkillToken economy:

\begin{enumerate}
  \item \textbf{Submit}: Validate the freelancer has sufficient token balance. Deduct tokens atomically. Create \texttt{Proposal} document in \texttt{PENDING} status. Create a \texttt{TokenTransaction} record.
  \item \textbf{Withdraw}: Only allowed if \texttt{status === PENDING}. Refund the original \texttt{tokenCost}. Update proposal status to \texttt{WITHDRAWN}. Create a \texttt{CREDIT} \texttt{TokenTransaction}.
  \item \textbf{Accept}: Transition proposal to \texttt{ACCEPTED}, create a \texttt{Contract}, transition project to \texttt{IN\_PROGRESS}, reject all other proposals.
\end{enumerate}

\section{Middleware Chain}

\subsection{Authentication Middleware}

\begin{lstlisting}[language=TypeScript, caption={JWT Authentication Middleware (simplified)}]
export const protect = async (req, res, next) => {
  const token = req.headers.authorization?.split(' ')[1];
  if (!token) return res.status(401).json({ message: 'Unauthorised' });
  const decoded = verifyAccessToken(token);
  req.user = decoded;
  next();
};
\end{lstlisting}

\subsection{Role Middleware}

\begin{lstlisting}[language=TypeScript, caption={Role Guard Middleware}]
export const requireRole = (role: Role) => (req, res, next) => {
  if (req.user?.role !== role)
    return res.status(403).json({ message: 'Forbidden' });
  next();
};
\end{lstlisting}

\subsection{Upload Middleware}

The \texttt{upload.middleware.ts} configures Multer with in-memory storage (\texttt{memoryStorage()}) and a 10MB file size limit. Files are stored as \texttt{Buffer} objects on \texttt{req.file}, which are then piped to Cloudinary via \texttt{streamifier.createReadStream()}.

\subsection{Request Logger}

The custom \texttt{requestLogger} middleware logs the HTTP method, URL, response status code, and response time in milliseconds to stdout for every request, aiding development debugging.

\section{Backend Directory Structure}

\begin{lstlisting}[caption={Backend Directory Structure (ASCII)}]
src/
  app.ts               -- Express app: middleware + route registration
  server.ts            -- HTTP server bootstrap + port binding
  config/
    cloudinary.ts      -- Cloudinary SDK init
    passport.ts        -- Google OAuth strategy definition
    prisma.ts          -- Prisma client singleton
    redis.ts           -- Redis client configuration
  controllers/         -- Route handler functions (thin layer)
    auth.controller.ts
    freelancer-onboarding.controller.ts
    freelancer.controller.ts
    project.controller.ts
    proposal.controller.ts
    token.controller.ts
    tracking.controller.ts
    contract.controller.ts
    metadata.controller.ts
  routes/              -- Express Router definitions
    auth.routes.ts
    freelancer.routes.ts
    project.routes.ts
    proposal.routes.ts
    token.routes.ts
    tracking.routes.ts
    category.routes.ts
    metadata.routes.ts
    ai.routes.ts
  middlewares/
    auth.middleware.ts      -- JWT verification
    role.ts                 -- Role-based access guard
    upload.middleware.ts    -- Multer file handling
    logger.middleware.ts    -- Request logging
  services/
    token.service.ts
    tracking.service.ts
  modules/
    browse/             -- Browse & discovery module
  utils/
    email.ts            -- Email templates
    jwt.ts              -- Token signing/verification helpers
    redis.ts            -- Redis helper wrappers
    uploadToCloudinary.ts
    profileCompletion.ts
    tokenCalculator.ts
    validators.ts
  ai/                   -- AI assistant module
\end{lstlisting}

\section{Utility Modules}

\begin{table}[H]
\centering
\caption{Utility Module Responsibilities}
\begin{tabularx}{\textwidth}{lX}
\toprule
\textbf{Utility} & \textbf{Responsibility} \\
\midrule
\texttt{jwt.ts}               & Wraps \texttt{jsonwebtoken} sign/verify. Access token: 15 min. Refresh token: 7 days. \\
\texttt{email.ts}             & HTML email templates for OTP, password reset, welcome, and notification emails. \\
\texttt{redis.ts}             & Helper wrappers: \texttt{setEx}, \texttt{get}, \texttt{del} for Redis OTP and cache storage. \\
\texttt{uploadToCloudinary.ts} & Uploads a Buffer to a specified Cloudinary folder and returns the secure URL. \\
\texttt{profileCompletion.ts} & Counts non-null mandatory fields to compute the 0--100 completion score. \\
\texttt{tokenCalculator.ts}   & Determines the SkillToken cost for a proposal based on project budget tiers. \\
\texttt{validators.ts}        & Zod schemas for request body validation across multiple controllers. \\
\bottomrule
\end{tabularx}
\end{table}

\section{Browse Module}

The \texttt{modules/browse/} directory contains a self-contained browse and discovery feature. It implements:

\begin{itemize}
  \item \textbf{Freelancer search}: Filter by skills, experience level, availability, minimum rating, hourly rate range, and location. Results are sorted by a composite \textit{relevance score} incorporating profile completeness, rating, dispute ratio, and recent login activity.
  \item \textbf{Project browse}: Filter by category, budget range, project size, experience level, and skills. Records are sorted by recency and proposal count.
  \item \textbf{Interaction tracking}: \texttt{BrowseInteraction} records are created when a freelancer views a project, feeding back into personalised recommendations.
\end{itemize}

\section{Cron Jobs}

A monthly cron job (scheduled via \texttt{node-cron}) awards SkillTokens to qualifying active freelancers. A freelancer qualifies if they have been active (\texttt{lastLoginAt} within the last 30 days) and have a validated profile. The job creates a \texttt{MONTHLY\_AWARD} \texttt{TokenTransaction} for each qualifying freelancer and updates their \texttt{skillTokenBalance}.
